\documentclass[11pt,a4paper]{article}
\usepackage[utf8]{inputenc}
\usepackage[T1]{fontenc}
\usepackage[margin=0.7in]{geometry}
\usepackage{enumitem}
\usepackage{titlesec}
\usepackage{xcolor}
\usepackage{hyperref}
\definecolor{primarycolor}{RGB}{33,37,41}
\definecolor{secondarycolor}{RGB}{0,92,153}
\definecolor{accentcolor}{RGB}{90,100,110}
\hypersetup{colorlinks=true,urlcolor=secondarycolor}
\titleformat{\section}{\large\bfseries\color{primarycolor}}{}{0em}{}
\titlespacing*{\section}{0pt}{0.9em}{0.35em}
\setlist[itemize]{leftmargin=1.4em,itemsep=0.15em,topsep=0.15em}
\pagestyle{empty}
\setlength{\parindent}{0pt}
\begin{document}
\begin{center}
{\LARGE\bfseries MARTA KOWAL-BRANDT}\\[0.25em]
{\large QA AUTOMATION ENGINEER | CI/CD AND PLATFORM INFRASTRUCTURE}\\[0.25em]
marta.kowalbrandt@example.com | +48 555 303 303 | gitlab.com/mkb-example | github.com/mkb-dev
\end{center}
\section*{PROFESSIONAL SUMMARY}
QA automation engineer with 7 years of experience, including the last two focused mostly on CI/CD pipelines, test environments and build tooling. Owns GitLab CI pipelines and builds Helm-based ephemeral environments within a shared Kubernetes cluster. Personal infrastructure work includes k3s, GitOps, Terraform and observability tooling. Targeting a DevOps or platform engineering role.
\section*{SKILLS}
\textbf{CI/CD and containers:} GitLab CI, Jenkins, Docker, Docker Compose\\
\textbf{Infrastructure and orchestration:} Terraform (personal projects), Helm, Kubernetes (namespace-level access and test-environment deployments), k3s, Argo CD, Ansible (basic)\\
\textbf{Monitoring:} Prometheus, Grafana, Loki\\
\textbf{Programming and data:} Python, Bash, Java (test code), SQL, PostgreSQL\\
\textbf{Testing:} pytest, Playwright, Selenium, Postman\\
\textbf{Systems and tools:} Linux, Git\\
\textbf{Cloud learning:} AWS free-tier learning, including EC2 and S3 basics
\section*{SELECTED INFRASTRUCTURE PROJECTS}
\textbf{Merge-Request Ephemeral Environments | Odrava Health Software | 2023–2024}
\begin{itemize}
\item Built short-lived environments for merge requests, with an application, PostgreSQL and seeded data, deployed via Helm to a shared Kubernetes cluster from GitLab CI.
\item Configured automatic teardown after merge or 48 hours of inactivity; QA and product teams use the environments for feature reviews and e2e tests.
\end{itemize}
\textbf{CI Pipeline Speed-up | Odrava Health Software | 2022}
\begin{itemize}
\item Profiled pipelines for three teams, parallelised test jobs, added dependency caching and split e2e tests into a separate stage.
\item Reduced average pipeline time from 38 to 14 minutes and monthly CI runner cost by about 30\%.
\end{itemize}
\textbf{Homelab k3s Cluster with GitOps | Personal | 2024–2025}
\begin{itemize}
\item Runs a three-node k3s cluster on Raspberry Pi 5 devices, with applications deployed from a Git repository using Argo CD.
\item Set up Prometheus, Grafana and Loki, plus Traefik ingress and cert-manager for TLS.
\item Provisioned infrastructure for a small Hetzner VPS with Terraform; used Ansible at a basic level.
\end{itemize}
\textbf{Flaky Test Dashboard | Odrava Health Software | 2023}
\begin{itemize}
\item Collected CI test results in PostgreSQL and built Grafana dashboards showing flaky tests by team.
\item Flaky test rate fell from 7\% to below 2\% after teams began fixing the leading offenders.
\end{itemize}
\section*{EXPERIENCE}
\textbf{QA Automation Engineer (Test Infrastructure) | Odrava Health Software | 2021–Present}
\begin{itemize}
\item Own GitLab CI pipelines for three product teams (about 25 developers), covering build, unit, integration and e2e test stages, and artifact publishing.
\item Build and maintain Helm-based ephemeral test environments using the company’s shared Kubernetes cluster; access is at namespace level.
\item Maintain a pytest and Playwright e2e suite of about 900 tests; reduced the flaky rate from 7\% to below 2\%.
\item Write Bash and Python tooling for PostgreSQL test-data seeding.
\item Wrote the internal “How our CI works” documentation, used in team onboarding.
\item Mentored two QA interns in 2023.
\end{itemize}
\textbf{QA Engineer | Brightmill Software House | 2019–2021}
\begin{itemize}
\item Tested web and mobile applications for client projects, using both manual and automated testing.
\item Introduced a Selenium and Java automated regression suite for the largest client project.
\item Configured Jenkins jobs for nightly regression runs.
\end{itemize}
\textbf{Customer Support Specialist (2nd line) | Brightmill Software House | 2018–2019}
\begin{itemize}
\item Investigated customer issues in logs and databases, and reproduced bugs for developers.
\end{itemize}
\section*{EDUCATION AND CERTIFICATIONS}
BEng, Computer Engineering | Silesian Lowlands University of Technology | 2018\\
ISTQB Foundation Level | 2019\\
Certified Kubernetes Administrator (CKA): studying; exam booked for January 2027
\section*{LANGUAGES}
Polish (native), English (C1), German (A2)
\end{document}